\problema{Pacific Atlantic Water Flow}{417}{src/01-grafos/417-pacific-atlantic.cpp}{M}

Nos dan una matriz \texttt{heights} de tamaño $m\times n$ con la altura de
cada celda. El océano Pacífico toca el borde superior y el borde izquierdo de
la matriz; el océano Atlántico toca el borde inferior y el borde derecho. El
agua fluye de una celda a una vecina (arriba, abajo, izquierda, derecha) si la
altura de la vecina es \textbf{menor o igual} a la de la celda actual (el
agua desciende o se mantiene, nunca sube). Queremos las coordenadas de las
celdas desde las que el agua puede llegar a \textbf{ambos} océanos.

\paragraph{Intuición.} Hacer un BFS/DFS de escurrimiento desde cada una de las
$m\cdot n$ celdas para ver a qué océano llega sería $O((m\cdot n)^2)$: muy
lento. La clave es invertir el problema: en lugar de preguntar ``¿a dónde
llega el agua desde aquí?'', partimos de los bordes de cada océano y
exploramos hacia adentro \emph{subiendo} de altura. Una celda vecina es
alcanzable en este recorrido inverso si su altura es mayor o igual a la de la
celda actual, porque en el sentido original el agua fluiría de esa vecina
hacia la celda que estamos visitando. Así obtenemos, en una sola pasada por
océano, el conjunto de celdas que logran drenar hasta él. La respuesta final
es la intersección de ambos conjuntos.

\begin{center}
\ttfamily
\begin{tabular}{ccccc}
1 & 2 & 2 & 3 & 5\\
3 & 2 & 3 & 4 & 4\\
2 & 4 & 5 & 3 & 1\\
6 & 7 & 1 & 4 & 5\\
5 & 1 & 1 & 2 & 4\\
\end{tabular}
\end{center}

\noindent En este ejemplo, el resultado es
$(0,4), (1,3), (1,4), (2,2), (3,0), (3,1), (4,0)$: desde cada una de estas
celdas el agua puede descender hasta tocar tanto el borde superior/izquierdo
como el inferior/derecho.

\paragraph{Diagrama.} DFS inverso desde cada océano, subiendo de altura hacia
adentro de la matriz:

\begin{center}
\ttfamily
\begin{tabular}{ll}
Pacífico (fila 0 y columna 0) & Atlántico (fila $m$-1 y columna $n$-1)\\
P P P P P & \phantom{A A A A} A\\
P \phantom{P P P} & \phantom{A A A A} A\\
P \phantom{P P P} & (subiendo de altura) \phantom{AA} A\\
P \phantom{P P P} & \phantom{A A A A} A\\
P \phantom{P P P} & A A A A A A A A\\
\end{tabular}
\end{center}

\paragraph{Solución.} Dos matrices booleanas \texttt{pacifico} y
\texttt{atlantico} marcan qué celdas alcanzan cada océano. Lanzamos un DFS
desde cada celda del borde correspondiente; el DFS sólo avanza hacia una
vecina cuya altura sea mayor o igual (nunca hacia una más baja, porque el
agua no podría subir en el sentido inverso). Al final, las celdas marcadas en
ambas matrices forman la respuesta:

\lstinputlisting{src/01-grafos/417-pacific-atlantic.cpp}

\complejidad{$O(m\cdot n)$}{$O(m\cdot n)$}

\paragraph{Notas de entrevista (Amazon).} El truco central es reconocer que
invertir el sentido del recorrido (partir de los océanos en vez de partir de
cada celda) baja la complejidad de $O((m\cdot n)^2)$ a $O(m\cdot n)$; es un
patrón que vale la pena mencionar en voz alta al entrevistador. Casos borde
importantes: matriz $1\times1$ (la única celda toca los cuatro bordes a la
vez), una sola fila o columna (la misma fila/columna es simultáneamente
borde de ambos océanos), y mesetas de altura uniforme (el agua fluye
libremente entre vecinas de igual altura porque la condición es ``menor o
igual'').
